% !TEX root = ../main.tex
\chapter{Learning-Rate Schedulers}
\label{ch:lr_schedulers}

\todo[inline]{Chapter 40 --- Learning-Rate Schedulers. Multi-level scaffold below; prose to follow.}

\section{Overview}
\label{sec:lr_schedulers:overview}

\section{Architecture}
\label{sec:lr_schedulers:arch}

\subsection{Core building blocks}
\label{ssec:lr_schedulers:blocks}

\subsection{Variants and parameterizations}
\label{ssec:lr_schedulers:variants}

\section{Forward Computation}
\label{sec:lr_schedulers:forward}

\subsection{Mathematical formulation}
\label{ssec:lr_schedulers:formulation}

\subsection{Implementation in Lucid}
\label{ssec:lr_schedulers:impl}

\subsubsection{Module class layout}
\subsubsection{Backend dispatch}

\section{Training Considerations}
\label{sec:lr_schedulers:training}

\subsection{Initialization strategy}
\label{ssec:lr_schedulers:init}

\subsection{Backward pass}
\label{ssec:lr_schedulers:backward}

\subsection{Interaction with optimizer and AMP}
\label{ssec:lr_schedulers:optim_amp}

\section{Implementation Notes}
\label{sec:lr_schedulers:notes}

\subsection{File layout and class hierarchy}
\label{ssec:lr_schedulers:layout}

\subsection{Edge cases and pitfalls}
\label{ssec:lr_schedulers:pitfalls}

\section{Worked Examples and Benchmarks}
\label{sec:lr_schedulers:examples}

\section{Summary and Cross-References}
\label{sec:lr_schedulers:summary}
